\section*{Přínos}
\addcontentsline{toc}{section}{Přínos}

Navrhované aplikační řešení se zaměřuje především na využití architektury k získávání informací o kvalitě prostředí, ve kterém se uživatel pohybuje. Tato architektura umožňuje shromažďovat a analyzovat data o kvalitě vzduchu a dalších faktorů, což poskytuje uživatelům přehled o aktuálním stavu jejich prostředí. Uživatelé mohou na základě těchto informací podniknout kroky k nápravě, pokud prostředí není vyhovující.

Díky otevřenosti dat o IEQ (Indoor Environmental Quality) může aplikace také sloužit jako nástroj pro výrobce topných, ventilačních a klimatizačních systémů (HVAC), kteří mohou na základě těchto reálných dat vylepšovat své produkty a přizpůsobit je konkrétním potřebám zákazníků. Tímto způsobem přispívá k vyšší efektivitě a komfortu pracovního a domácího prostředí, což může vést ke zlepšení pracovní výkonosti a pohodlí nejen zaměstnanců, ale i studentů a dalších uživatelů.

Tato architektura tedy není jenom prostředkem pro implementaci aplikace, ale klíčovým nástrojem pro pochopení a zlepšení kvality vnitřního prostředí, čímž podporuje vyšší úroveň komfortu a zdraví uživatelů v každodenním životě.